\section{Limitations and Future Work}
\label{sec:limitations-future-work}
\ours{} represents our pioneering effort on facilitating mixed-initiative human-robot interaction through mixed-initiative natural-language dialog. While we focused on delegating steps for long-horizon manipulation tasks in a manner that maximizes task success and minimizes human effort, we believe this paper opens up exciting new avenues for future work. These include enabling both agents to learn to provide and incorporate spatial-temporal feedback to each other while performing a task, share relevant task information in an imperfect-information setting, and replan and redefine a task as necessary, all through mixed-initiative dialog.

\ours{} has a number of limitations.
% First, it assumes a fixed plan with a predetermined ordering of steps, where the human has a general, high-level understanding of the plan (but not the low-level steps that the robot plans over). We leave replanning, such as to react or better accommodate ad-hoc human requests, to future work.
% It cannot handle cases where the human wishes to add or remove steps from the plan dynamically, such as if the user tells the robot to ``grab another cold drink while you're at the fridge before coming back to me.''
First, it assumes that the human and robot work sequentially, and cannot handle cases where a robot and human wish to collaborate simultaneously on the same step in the plan, such as if the robot holds a roll of tape and the human cuts from it. 
Second, \ours{} assumes a fixed plan with a predetermined ordering of steps, where the human has a general, high-level understanding of the plan (but not the low-level steps that the robot plans over).
% Furthermore, \ours{} does not support parallelization where both the human and robot can work on different steps of a task simultaneously. One way to address this would be to operate on plan trees, where the parent nodes are steps that must be done before the child nodes, and sibling nodes can be executed by either agent in parallel.
Our method could be improved with a more nuanced definition of ``effort'' beyond our time-based metric.
% For instance, \ours{} assumes that ``effort'' is based on the time necessary to perform a task until completion. However, effort may also depend on the intensity of the task, how much the user enjoys it, and how physically capable each user is---our method had sidestepped this issue by assuming each human would expend the same amount of effort for each action primitive from some given state $s$.
Finally, $p_{H, t}$ prediction can be improved, such as by processing tone-of-voice and facial expressions, to enable producing more emotionally understanding dialog, which can improve task success and user satisfaction.